%------------------------------------------------------------------------------%
\section{Revisão}

\begin{exercises*}

% Thomas pg 72 1-18
%------------------------------------------------------------------------------%
\begin{exercise}
  Determine quais sequências são convergentes e quais são divergentes.
  Calcule o limite das sequências convergentes.
  \begin{mathlist}
    % 1 - 4
    \item a_n = 1 + \frac{(-1)^n}{n}
    \item a_n = \frac{1-(-1)^n}{\sqrt{n}}
    \item a_n = \frac{1-2^n}{2^n}
    \item a_n = 1+\left(\frac{9}{10}\right)^n
    % 5-8
    \item a_n = \sin\left(\frac{n\pi}{2}\right)
    \item a_n = \sin(n\pi)
    \item a_n = \frac{\ln\left(n^2\right)}{n}
    \item a_n = \frac{\ln(2n+1)}{n}
    % 9 - 12
    \item a_n = \frac{n+\ln(n)}{n}
    \item a_n = \frac{\ln\left(2n^3+1\right)}{n}
    \item a_n = \left(\frac{n-5}{n}\right)^n
    \item a_n = \left(1+\frac{1}{n}\right)^{-n}
    % 13-15
    \item a_n = \sqrt[n]{\frac{3^n}{n}}
    \item a_n = \left(\frac{3}{n}\right)^{\frac{1}{n}}
    \item a_n = n \left(2^{\sfrac{1}{n}}-1\right)
    % 16 - 18
    \item a_n = \sqrt[n]{2n+1}
    \item a_n = \frac{(n+1)!}{n!}
    \item a_n = \frac{(-4)^n}{n!}
  \end{mathlist}
\end{exercise}

% Exercício de prova do prof. André Pereira
%------------------------------------------------------------------------------%
\begin{exercise}
Considere as sequências
\[
  a_n = \ln\left(1+\frac{1}{n^2}\right)
  \qquad\text{e}\qquad
  b_n = \frac{1}{n^2}
\]
\begin{alphalist}
  \item Enuncie a definição de convergência de uma sequência e use essa definição
        para mostrar que $b_n\rightarrow 0$. \\
        \hint{Se $\varepsilon>0$ e $n>1/\sqrt{\varepsilon}$ então $\sfrac{1}{n^2}<\varepsilon$.}
  \item Mostre que $a_{n}\rightarrow 0$ e que \(\frac{a_n}{b_n}\rightarrow 1\).
\end{alphalist}
\end{exercise}

% Exercício de prova do prof. André Pereira
%------------------------------------------------------------------------------%
\begin{exercise}
    Defina divergência ao infinito (\(\lim a_n=\infty\)).
    Mostre, usando essa definição, que se $\lim a_n=\infty,$ então
    \(\lim \frac{1}{a_n}=0\).
\end{exercise}

% Exercício de prova do prof. André Pereira
%------------------------------------------------------------------------------%
\begin{exercise}
    Considere as sequências
    \[
        a_n = \frac{n!}{5^n}    \hfill
        b_n = \frac{1}{n\ln n}  \hfill
        c_n = \left(\frac{1}{n}\right)^{1/(\ln n)} \hfill
    \]
    Verifique se as sequências $a_n,$ $b_n$ e $c_n$ convergem ou divergem.
    Calcule o limite no(s) caso(s) em que converge.
\end{exercise}

% Exercício de prova do prof. André Pereira
%------------------------------------------------------------------------------%
\begin{exercise}
  Expresse a dizima periódica $0,\overline{23}=0,232\,323\,232\dots$
  como a razão de dois números inteiros.
\end{exercise}

% Exercício de prova do prof. André Pereira
%------------------------------------------------------------------------------%
\begin{exercise}
  Qual frase a seguir é verdadeira e qual é falsa?
  Prove a verdadeira e dê um contra exemplo para a falsa.
  \begin{alphalist}
    \item Toda sequência convergente é limitada.
    \item Toda sequência limitada é convergente.
  \end{alphalist}
\end{exercise}

% Exercício de prova do prof. André Pereira
%------------------------------------------------------------------------------%
\begin{exercise}
Considerando as sequências
\[
  a_n = \frac{(n!)^2}{(2n)!}            \hfill
  b_n = \frac{1}{\sqrt{n}}              \hfill
  c_n = \frac{1}{2\sqrt{n}+\sqrt[3]{n}} \hfill
\]
mostre que
\[
  a_n\rightarrow 0                        \hfill
  b_n^{\sfrac{2}{n}}\rightarrow 1         \hfill
  \frac{c_n}{b_n}\rightarrow \frac{1}{2}  \hfill
\]
\end{exercise}

% Exercício de prova do prof. André Pereira
%------------------------------------------------------------------------------%
\begin{exercise}
    Mostre que $1=0,\bar{9}=0,999\dots$
\end{exercise}

% 44-56 - Stewart pt pg 633
%------------------------------------------------------------------------------%
\begin{exercise}
  Verifique se a sequência converge ou não e calcule o limite das sequências convergentes.
  \begin{mathlist}
    \item a_n = \sqrt[n]{2^{1+3n}\,}
    \item a_n = n\sin\left(\frac{1}{n}\right)
    \item a_n = \left( 1+ \frac{2}{n}\right)^n
    \item a_n = \ln\left(2n^2+1\right) - \ln\left(n^2+1\right)
    \item a_n = \frac{\big(\ln(n)\big)^2}{n}
    \item a_n = \arctan\big(\ln(n)\big)
    \item a_n = n - \sqrt{n+1\,}\sqrt{n+3\,}
    % 55
    \item a_n = \frac{n!}{2^n}
    \item a_n = \frac{(-3)^n}{n!}
  \end{mathlist}
\end{exercise}


%------------------------------------------------------------------------------%
%\begin{exercise}
%\begin{answer}
%\end{answer}
%\end{exercise}

%------------------------------------------------------------------------------%
%\begin{exercise}
%\begin{mathlist}[2]
%  \item
%  \item
%  \item
%\end{mathlist}
%\begin{answer}
%  \begin{alphalist}[3]
%    \item
%    \item
%    \item
%  \end{alphalist}
%\end{answer}
%\end{exercise}

\end{exercises*}
%------------------------------------------------------------------------------%
